\documentclass[11pt,a4paper]{article}
\usepackage[utf8]{inputenc}
\usepackage[english]{babel}
\usepackage{amsmath, amssymb, amsfonts}
\usepackage{graphicx}
\usepackage{booktabs}
\usepackage{float}
\usepackage{geometry}
\usepackage{hyperref}
\usepackage{microtype}

\geometry{top=2.5cm, bottom=2.5cm, left=2cm, right=2cm}
\hypersetup{
    colorlinks=true,
    linkcolor=blue,
    citecolor=blue,
    urlcolor=blue
}

\title{\textbf{Modelling and Forecasting the Norwegian Stock Index: A Comparative Approach between SARIMAX, VAR and Recurrent Neural Networks (LSTM)}}

\author{
  \textbf{André Marques} \\
  Master's in Data Science \\
  University Institute of Lisbon (ISCTE-IUL) \\
  \texttt{diana.mendes@iscte-iul.pt}
}

\date{31 May 2026}

\begin{document}

\maketitle

\begin{abstract}
This work presents an empirical study on the modelling and forecasting of Norway's stock index (\textit{Share Price Index}), integrating its macro-financial relationship with the price of Brent crude oil and the USD/NOK exchange rate over the period from 1992 to 2026. As Norway is one of Europe's main oil-exporting economies, this study analyses the dynamic interactions between these variables. We employ three distinct classes of models: a univariate model with exogenous variables (SARIMAX), a classical linear multivariate model (Vector Autoregression - VAR) and a non-linear \textit{Deep Learning} model (LSTM recurrent network in PyTorch). Our out-of-sample results show the clear superiority of multivariate approaches in the presence of macro-financial dynamics, with the LSTM network achieving the highest forecast accuracy with a MAPE of only 2.62\%, closely followed by the VAR model (2.94\%).
\end{abstract}

\section{Introduction}
Norway's capital market has unique characteristics owing to the country's economic profile. Norway is one of the largest oil and gas exporters in the world, which makes its currency (the Norwegian Krone - NOK) a ``petrocurrency'' and intrinsically links the performance of its listed companies (represented in the Share Price Index) to fluctuations in energy \textit{commodities}.

In this context, we pose the following \textbf{research questions}:
\begin{enumerate}
    \item \textit{To what extent do the price of Brent crude oil and the USD/NOK exchange rate influence and help forecast the Norwegian stock index?}
    \item \textit{Can recurrent neural networks (LSTM) capture complex non-linear relationships better than classical econometric models (VAR and SARIMAX) in this market?}
\end{enumerate}

The main objective of this work is to model and forecast the monthly time series of the Norwegian Share Price Index in a robust way. The main contributions of this study are: (i) to show empirically the strong short-term link between oil and the Norwegian financial economy, (ii) to demonstrate that multivariate modelling in first differences of logs is essential to stabilise the variance and forecast macroeconomic trends, and (iii) to present a rigorous out-of-sample comparison between statistical models and \textit{Deep Learning} algorithms.

The rest of the article is structured as follows: Section 2 presents the theoretical and empirical background of the models; Section 3 details the architecture of the fitted models; Section 4 describes the experimental setup and data collection; Section 5 analyses the results and diagnostics in detail; and Section 6 concludes the study, presenting its limitations and future directions.

\section{Theoretical Background}
To model the Norwegian stock index ($S_t$), the price of Brent crude oil ($O_t$) and the USD/NOK exchange rate ($E_t$), we use three essential methodologies:

\subsection{SARIMAX Model}
The Seasonal Autoregressive Integrated Moving Average model with Exogenous variables (SARIMAX) is a univariate extension of the classical Box-Jenkins ARIMA models. Mathematically, for a series $y_t$:
\begin{equation}
\phi_p(L)\Phi_P(L^s)(1-L)^d(1-L^s)^D y_t = \beta X_t + \theta_q(L)\Theta_Q(L^s)\epsilon_t
\end{equation}
where $X_t$ represents the exogenous variable (in this case, the log of Brent crude oil), $L$ is the lag operator ($L y_t = y_{t-1}$), and $\epsilon_t$ represents white noise. This model assumes linearity and focuses on the autocorrelation of the variable of interest itself, adding the impact of external factors.

\subsection{VAR Model (Vector Autoregression)}
The VAR model captures dynamic linear interdependencies between multiple time series. Unlike SARIMAX, all variables are treated as endogenous. A VAR with $p$ lags is expressed as:
\begin{equation}
\mathbf{Y}_t = \mathbf{c} + \mathbf{A}_1 \mathbf{Y}_{t-1} + \dots + \mathbf{A}_p \mathbf{Y}_{t-p} + \mathbf{u}_t
\end{equation}
where $\mathbf{Y}_t = [\Delta \log S_t, \Delta \log E_t, \Delta \log O_t]^T$, $\mathbf{A}_i$ are coefficient matrices to be estimated, and $\mathbf{u}_t$ is a vector of zero-mean innovations. If there were cointegration between the variables in levels, the correct specification would be the VECM (Vector Error Correction Model), which includes a long-run error correction term. However, in the absence of cointegration, the VAR in differences is the theoretically appropriate specification to avoid spurious regressions.

\subsection{Recurrent Neural Networks (LSTM)}
\textit{Long Short-Term Memory} (LSTM) networks are a \textit{Deep Learning} architecture designed specifically for sequential data, able to mitigate the vanishing gradient problem. The mathematical intuition is based on the use of dynamic \textit{gates}: the forget gate ($f_t$), the input gate ($i_t$) and the output gate ($o_t$). These control the flow of information in the cell state ($c_t$), allowing the network to retain non-linear long-term dependencies, which are crucial for the volatility of financial markets:
\begin{equation}
f_t = \sigma(W_f [h_{t-1}, x_t] + b_f)
\end{equation}

\section{Models and Architecture}
The final models were fitted on the training set following rigorous optimisation criteria:

\begin{enumerate}
    \item \textbf{SARIMAX(1,1,1)x(0,0,0)[12]}: Determined automatically by minimising the AIC through the Stepwise algorithm. It uses $\log(O_t)$ as an exogenous regressor. The mathematical specification shows the need for a first difference ($d=1$) to stabilise the mean.
    \item \textbf{VAR(2)}: Lag selection was based on the AIC information criterion, identifying $p=2$ monthly lags as the optimal structure to capture the momentum effect of oil and the exchange rate on the Norwegian stock market.
    \item \textbf{LSTM Forecaster}: Implemented in PyTorch, the final architecture consists of an LSTM layer with $8$ hidden neurons (\textit{hidden dimension}), a \textit{lookback} window of 3 months, and a linear dense output layer for forecasting the log-return. The loss function was the MSE (Mean Squared Error), optimised with Adam at a learning rate of $\alpha = 0.01$ over 150 epochs.
\end{enumerate}

\section{Experimental Setup}
\subsection{Data Collection and Alignment}
We downloaded 3 monthly series from FRED for the period from 1 January 1992 to 1 March 2026 (411 complete observations):
\begin{itemize}
    \item \textbf{Share\_Index} ($S_t$): SPASTT01NOM661N (OECD Stock Price Index)
    \item \textbf{USD\_NOK} ($E_t$): EXNOUS (USD/NOK exchange rate)
    \item \textbf{Brent\_Oil} ($O_t$): POILBREUSDM (Global Price of Brent Crude)
\end{itemize}

\subsection{Descriptive Statistics}
Table \ref{tab:desc} summarises the descriptive statistics of the training series (369 observations, 90\% of the sample). The stock index shows moderate positive skewness ($0.92$), while the exchange rate and oil have negative kurtosis, reflecting marked economic cycles.

\begin{table}[ht]
\centering
\small
\begin{tabular}{lrrrrr}
\toprule
\textbf{Variable} & \textbf{Mean} & \textbf{Std. Dev.} & \textbf{Minimum} & \textbf{Maximum} & \textbf{Skewness} \\
\midrule
$S_t$ (Stocks) & 66.80 & 52.02 & 6.73 & 216.43 & 0.92 \\
$E_t$ (Exchange rate) & 7.23 & 1.20 & 5.05 & 10.42 & 0.35 \\
$O_t$ (Oil) & 53.01 & 32.96 & 10.16 & 133.59 & 0.57 \\
\bottomrule
\end{tabular}
\caption{Descriptive Statistics of the Variables in Levels (Training).}
\label{tab:desc}
\end{table}

\subsection{Stationarity Tests}
The Augmented Dickey-Fuller (ADF), Phillips-Perron (PP) and KPSS tests showed empirically that all three series in levels are \textbf{non-stationary} (Table \ref{tab:roots}). After applying the first difference of the natural logarithm ($\Delta \log$), all series reject the unit root test at 1\%, becoming stationary ($I(1)$).

\begin{table}[ht]
\centering
\small
\begin{tabular}{lrrr}
\toprule
\textbf{Series (Training)} & \textbf{ADF (p-val)} & \textbf{PP (p-val)} & \textbf{KPSS (Stat)} \\
\midrule
$S_t$ (Level) & 0.8481 & 0.8815 & 0.4924* \\
$E_t$ (Level) & 0.9060 & 0.8749 & 0.4298* \\
$O_t$ (Level) & 0.1350 & 0.3787 & 0.3471* \\
\midrule
$\Delta \log S_t$ (Diff) & 0.0000 & 0.0000 & 0.0291 \\
$\Delta \log E_t$ (Diff) & 0.0000 & 0.0000 & 0.0855 \\
$\Delta \log O_t$ (Diff) & 0.0000 & 0.0000 & 0.0451 \\
\bottomrule
\end{tabular}
\caption{Unit Root Test Results (Level vs. 1st Difference of the Log). * indicates rejection of stationarity at 5\%.}
\label{tab:roots}
\end{table}

\subsection{Cointegration Test}
The Johansen cointegration test applied to the vector $[\log S_t, \log E_t, \log O_t]^T$ produced a trace statistic of $15.30$ for rank $r \leq 0$, well below the 95\% critical value of $29.80$. We therefore conclude that there is \textbf{no cointegration}, which theoretically justifies estimating the VAR model in first differences.

\section{Results and Analysis}
The test (\textit{out-of-sample}) period runs from October 2022 to March 2026 (42 observations). Table \ref{tab:errors} shows the comparison of the error metrics.

\begin{table}[ht]
\centering
\small
\begin{tabular}{lrrr}
\toprule
\textbf{Model} & \textbf{RMSE} & \textbf{MAE} & \textbf{MAPE (\%)} \\
\midrule
SARIMAX & 65.51 & 53.90 & 20.35\% \\
\textbf{VAR (2)} & 13.00 & 7.89 & 2.94\% \\
\textbf{LSTM (PyTorch)} & \textbf{11.29} & \textbf{6.87} & \textbf{2.62\%} \\
\bottomrule
\end{tabular}
\caption{Out-of-Sample Forecast Error Metrics in the Test Period.}
\label{tab:errors}
\end{table}

The analysis reveals a substantial gap between the univariate SARIMAX model and the multivariate models (VAR and LSTM). SARIMAX obtained a MAPE of $20.35\%$, failing to anticipate the structural changes and the post-pandemic and geopolitical accelerations in share prices. This is because SARIMAX assumes simplified linear dynamics for the exogenous variable.

By contrast, the multivariate \textbf{VAR(2)} model performed remarkably well, with a MAPE of only $2.94\%$, showing that capturing the cross-lags between Brent and the exchange rate is essential to model the Norwegian market. The \textbf{LSTM network} achieved the best absolute performance ($MAPE = 2.62\%$). The LSTM model's ability to retain long-term memory in its gates and to model non-linear interactions allowed it to track the actual values of the stock index very closely over the 2022-2026 period.

\section{Conclusion and Future Work}
This work showed that multivariate macro-financial modelling is highly effective for forecasting the Norwegian stock index, with the LSTM recurrent network and the VAR model standing out. The absence of cointegration indicates that these relationships lie mainly in short-term dynamics and percentage changes (\textit{log-returns}).

The main limitation of this work is its reliance on monthly frequency, which smooths out high-frequency peaks. For future research, we suggest: (i) modelling conditional volatility with a hybrid VAR-GARCH or LSTM-GARCH specification to capture volatility clusters, and (ii) extending the data collection to daily data, including a news sentiment index specific to the oil market.

\end{document}
